\documentclass[aspectratio=169]{beamer}

\setbeamertemplate{navigation symbols}{}
\usetheme{Boadilla}
\usepackage[italian]{babel}
\usepackage[dvipsnames]{xcolor}
\usepackage{minted}
\usepackage{calc} 

\title{Gestione Documenti - Versione RIA}
\author{Lorenzo Chiroli (C.P. 10797603 - MTR. 981392)}
\date{Gruppo N. 120}

\begin{document}
\frame{\titlepage}

\begin{frame}
    \frametitle{Specification and Optional Additions - 1}
    \textbf{Specifiche generali:}\\
    Per rendere la documentazione maggiormente concisa, si è deciso di non ripetere inutilmente le specifiche
    "generali" - già trattate nei lucidi relativi alla versione "Pure HTML" - e di concentrarsi solamente sulle
    differenze proprie della versione RIA.\\

    \textbf{Specifiche aggiuntive:}
    \begin{itemize}
        \item L'applicazione supporta registrazione e login mediante una pagina pubblica con opportune
              form. La registrazione controlla la validità sintattica dell'indirizzo di email e l'uguaglianza
              tra i campi "password" e "ripeti password", anche a lato client. La registrazione controlla
              l'unicità dello username.
        \item Dopo il login dell'utente, l'intera applicazione è realizzata con un'unica pagina.
        \item Ogni interazione dell'utente è gestita senza ricaricare completamente la pagina, ma produce
              l'invocazione asincrona del server e l'eventuale modifica del contenuto da aggiornare a
              seguito dell'evento.
    \end{itemize}
\end{frame}

\begin{frame}
    \frametitle{Specification and Optional Additions - 2}
    \begin{itemize}
        \item Errori a lato server devono essere segnalati mediante un messaggio di allerta all'interno
              della pagina.
        \item La funzione di creazione di una sottocartella è realizzata nella pagina HOME mediante un
              bottone AGGIUNGI SOTTOCARTELLA posto di fianco ad ogni cartella. La pressione del
              bottone fa apparire un campo di input per l'inserimento del nome della cartella da inserire.
        \item La funzione di creazione di un documento è realizzata nella pagina HOME mediante un
              bottone AGGIUNGI DOCUMENTO posto di fianco ad ogni cartella. La pressione del
              bottone fa apparire una form di input per l'inserimento dei dati del documento.
        \item Si aggiunge una cartella denominata "cestino". Il drag and drop di un documento o di una
              cartella nel cestino comporta la cancellazione. Prima di inviare il comando di cancellazione
              al server l'utente vede una finestra modale di conferma e può decidere se annullare
              l'operazione o procedere. La cancellazione di una cartella comporta la cancellazione
              integrale e ricorsiva del contenuto dalla base di dati (documenti e cartelle).
    \end{itemize}
\end{frame}

\begin{frame}
    \frametitle{Specification and Optional Additions - 3}
    \begin{itemize}
        \item La funzione di spostamento di un documento è realizzata mediante drag and drop.
    \end{itemize}
    \vspace{0.4cm}
    Per implementare questa funzionalità si è deciso di:
    \begin{enumerate}
        \item Nella homepage viene aggiunto un cestino in cui è possibile trascinare le cartelle per eliminarle.
        \item Nella vista "contenuto" di una cartella, quando l'utente inizia il trascinamento di un documento,
              viene visualizzata - in una "finestra modale" - la lista delle cartelle in cui è possibile spostare il
              documento (tra cui anche il cestino). La directory structure viene subito nascosta se l'utente annulla
              il drag\&drop rilasciando il documento.
    \end{enumerate}
\end{frame}

\begin{frame}
    \frametitle{Data Requirements Analysis - Database Design - Database Schema}
    \begin{itemize}
        \item La struttura della base di dati è identica a quella della versione "Pure HTML".
        \item A onore del vero, tutta la "buisness logic" server-side - Bean, DAO e utility varie - è stata mantentura
              inalterata e viene gestita come submodule comune a tutte le versioni del DMS per evitare ridondanze nella
              codebase.
    \end{itemize}
\end{frame}

\begin{frame}
    \frametitle{Application Requirements Analysis - 1}
    \begin{itemize}
        \item L'applicazione supporta \textcolor{Blue}{registrazione} e \textcolor{Blue}{login} mediante una
              \underline{pagina pubblica} [rif. a \textcolor{Red}{authentication.html}] con
              \textcolor{Green}{opportune form}. La registrazione \textcolor{Orange}{controlla la validità sintattica
                  dell'indirizzo di email e l'uguaglianza tra i campi "password" e "ripeti password", anche a lato client}.
              La registrazione \textcolor{Orange}{controlla l’unicità dello username}.
        \item Dopo il \textcolor{Blue}{login} dell'utente, l'intera applicazione è realizzata con un'
              \underline{unica pagina} [rif. a \textcolor{Red}{index.html}].
        \item Ogni interazione dell'utente è gestita senza ricaricare completamente la pagina, ma produce
              l'invocazione asincrona del server e l'eventuale modifica del contenuto da aggiornare a
              seguito dell'evento.
        \item Errori a lato server devono essere segnalati mediante un \textcolor{Green}{messaggio di allerta}
              all'interno della pagina.
        \item La funzione di \textcolor{Orange}{spostamento di un documento} è realizzata mediante
              \textcolor{Blue}{drag and drop}.
    \end{itemize}
    \footnotetext{I colori hanno il seguente significato: \textcolor{Red}{pagine},
        \textcolor{Green}{componenti di visualizzazione}, \textcolor{Blue}{eventi} e \textcolor{Orange}{azioni}.}
\end{frame}

\begin{frame}
    \frametitle{Application Requirements Analysis - 2}
    \begin{itemize}
        \item La funzione di \textcolor{Orange}{creazione di una sottocartella} è realizzata nella
              \underline{pagina HOME} [\textcolor{Red}{index.html}] mediante un
              \textcolor{Green}{bottone AGGIUNGI SOTTOCARTELLA} posto di fianco ad ogni cartella. La
              \textcolor{Blue}{pressione del bottone} \textcolor{Orange}{fa apparire} un
              \textcolor{Green}{campo di input} per l'inserimento del nome della cartella da inserire.
        \item La funzione di \textcolor{Orange}{creazione di un documento} è realizzata nella
              \underline{pagina HOME} [\textcolor{Red}{index.html}] mediante un
              \textcolor{Green}{bottone AGGIUNGI DOCUMENTO} posto di fianco ad ogni cartella.
              La \textcolor{Blue}{pressione del bottone} \textcolor{Orange}{fa apparire}
              una \textcolor{Green}{form di input} per l'inserimento dei dati del documento.
        \item Si aggiunge una \textcolor{Green}{cartella denominata "cestino"}. Il
              \textcolor{Blue}{drag and drop di un documento o di una cartella nel cestino} comporta
              \textcolor{Orange}{la cancellazione}. Prima di
              \textcolor{Orange}{inviare il comando di cancellazione al server} l'utente vede una
              \textcolor{Green}{finestra modale di conferma} e può
              \textcolor{Blue}{decidere se annullare l'operazione o procedere}. La cancellazione di una cartella
              comporta la cancellazione integrale e ricorsiva del contenuto dalla base di dati (documenti e cartelle).
    \end{itemize}
    \footnotetext{I colori hanno il seguente significato: \textcolor{Red}{pagine},
        \textcolor{Green}{componenti di visualizzazione}, \textcolor{Blue}{eventi} e \textcolor{Orange}{azioni}.}
\end{frame}

\begin{frame}
    \frametitle{Application Design - Components (1)}
    \textbf{Sommario dei componenti:}
    \begin{itemize}
        \item \textcolor{Red}{authentication.html}: Pagina di registrazione e login.
              \begin{itemize}
                  \item \textcolor{Green}{Form di Login}: Form per il login.
                  \item \textcolor{Green}{Bottone di Login}
                        \begin{itemize}
                            \item \textcolor{Blue}{(onClick)}: Esegue il login dell'utente.
                        \end{itemize}
                  \item \textcolor{Green}{Link di Registrazione}
                        \begin{itemize}
                            \item \textcolor{Blue}{(onClick)}: Mostra i componenti relativi alla registrazione
                                  nascondendo quelli del login.
                        \end{itemize}

                  \item \textcolor{Green}{Form di Registrazione}: Form per la registrazione.
                  \item \textcolor{Green}{Bottone di Registrazione}
                        \begin{itemize}
                            \item \textcolor{Blue}{(onClick)}: Esegue la registrazione dell'utente.
                        \end{itemize}
                  \item \textcolor{Green}{Link di Login}
                        \begin{itemize}
                            \item \textcolor{Blue}{(onClick)}: Mostra i componenti relativi al login
                                  nascondendo quelli di registrazione.
                        \end{itemize}
              \end{itemize}
    \end{itemize}
\end{frame}

\begin{frame}
    \frametitle{Application Design - Components (2)}
    \textbf{\textcolor{Red}{index.html}}: Pagina principale dell'applicazione.
    \begin{itemize}
        \item \textcolor{Green}{"Directory Structure"}: Albero delle cartelle dell'utente.
              \begin{itemize}
                  \item \textcolor{Blue}{(onClick)}: Mostra il contenuto (cartelle e documenti) della cartella
                        selezionata.
                  \item \textcolor{Blue}{(onDragStart)}: Inizia il trascinamento di una cartella o documento.
              \end{itemize}

        \item \textcolor{Purple}{Bottoni di Aggiunta}:
              \begin{itemize}
                  \item \textcolor{Green}{"Add New Folder"}: Bottone per aggiungere una nuova cartella di primo livello.
                        \begin{itemize}
                            \item \textcolor{Blue}{(onClick)}: Mostra un \textcolor{Green}{campo di input} per
                                  l'inserimento del nome della cartella da aggiungere.
                        \end{itemize}
                  \item \textcolor{Green}{"Add Subfolder"}: Bottone per aggiungere una sottocartella.
                        Uno per ogni cartella.
                        \begin{itemize}
                            \item \textcolor{Blue}{(onClick)}: Mostra un \textcolor{Green}{campo di input} per
                                  l'inserimento del nome della sottocartella da aggiungere.
                        \end{itemize}
                  \item \textcolor{Green}{"Add Document"}: Bottone per aggiungere un documento.
                        Uno per ogni cartella.
                        \begin{itemize}
                            \item \textcolor{Blue}{(onClick)}: Mostra una \textcolor{Green}{form di input} per
                                  l'inserimento dei dati del documento da aggiungere.
                        \end{itemize}
              \end{itemize}

        \item \textcolor{Green}{Cestino}: Cestino in cui trascinare le cartelle o i documenti per eliminarli.
              \begin{itemize}
                  \item \textcolor{Blue}{(onDrop)}: Mostra la finestra modale di conferma per la cancellazione
                        ("Confirm Deletion").
              \end{itemize}

        \item \textcolor{Green}{Finestra Modale - "Confirm Deletion"}: Finestra modale di conferma per
              la cancellazione.
              \begin{itemize}
                  \item \textcolor{Blue}{(onClick)}\textrightarrow \textcolor{Green}{Bottone "Confirm"}:
                        Elimina definitivamente l'elemento.
                  \item \textcolor{Blue}{(onClick)}\textrightarrow \textcolor{Green}{Bottone "Cancel"}:
                        Chiude la finestra modale senza effettuare la cancellazione.
              \end{itemize}
    \end{itemize}
\end{frame}

\begin{frame}
    \frametitle{Application Design - Components (3)}
    \begin{itemize}
        \item \textcolor{Green}{Lista Subfolder}: Lista delle sottocartelle di una cartella.
              \begin{itemize}
                  \item \textcolor{Blue}{(onClick)}: Mostra il contenuto (cartelle e documenti) della
                        sottocartella selezionata.
              \end{itemize}

        \item \textcolor{Green}{Lista Documenti}: Lista dei documenti di una cartella.
              \begin{itemize}
                  \item \textcolor{Blue}{(onClick)}: Mostra i dettagli del documento selezionato.
                  \item \textcolor{Blue}{(onDragStart)}: Mostra una finestra modale con la lista delle cartelle
                        in cui è possibile spostare il documento nonché il cestino.
              \end{itemize}

        \item \textcolor{Green}{Finestra Modale - "Move Document"}: Finestra modale per lo spostamento di un documento:
              contiene la lista delle cartelle dell'utente e il cestino.
              \begin{itemize}
                  \item \textcolor{Blue}{(onDrop)}\textrightarrow \textcolor{Green}{Subfolder}:
                        Mostra la finestra modale di conferma per lo spostamento del documento.
                  \item \textcolor{Blue}{(onDrop)}\textrightarrow \textcolor{Green}{Cestino}:
                        Mostra la finestra modale di conferma per la cancellazione del documento.
              \end{itemize}

        \item \textcolor{Green}{Finestra Modale - "Confirm Move"}: Finestra modale di conferma per lo spostamento
              di un documento.
              \begin{itemize}
                  \item \textcolor{Blue}{(onClick)}\textrightarrow \textcolor{Green}{Bottone "Confirm"}:
                        Sposta il documento nella cartella selezionata.
                  \item \textcolor{Blue}{(onClick)}\textrightarrow \textcolor{Green}{Bottone "Cancel"}:
                        Chiude la finestra modale senza effettuare lo spostamento.
              \end{itemize}
    \end{itemize}
\end{frame}

\begin{frame}
    \frametitle{Application Design - Components (4)}
    \begin{itemize}
        \item \textcolor{Green}{Form di Dettaglio Documento}: Form per la visualizzazione dei dettagli di un documento.

        \item \textcolor{Purple}{Bottoni di Navigazione}:
              \begin{itemize}
                  \item \textcolor{Green}{"Go Back"}: Bottone per tornare alla visualizzazione precedente.
                  \item \textcolor{Green}{"Home"}: Bottone per tornare alla visualizzazione della struttura cartelle.
                  \item \textcolor{Green}{"Logout"}: Bottone per effettuare il logout dall'applicazione.
              \end{itemize}
    \end{itemize}
\end{frame}

\end{document}